% Part: first-order-logic
% Chapter: axiomatic-deduction
% Section: deduction-theorem-quantifiers

\documentclass[../../../include/open-logic-section]{subfiles}

\begin{document}

\olfileid{fol}{axd}{ddq}

\olsection{带量词的演绎定理}

\begin{thm}[演绎定理]
\ollabel{thm:deduction-thm-q} 如果 $\Gamma \cup \{!A\} \Proves !B$，
则 $\Gamma \Proves !A \lif !B$。
\end{thm}

\begin{proof}
我们再次对从 $\Gamma \cup \{!A\}$ 推出 $!B$ 的推导长度进行归纳。

归纳基始的证明与 \olref[ded]{thm:deduction-thm} 中的证明相同。

对于归纳步骤，同样假设从 $\Gamma \cup \{!A\}$ 推出 $!B$ 的推导的最后一步~$!B$ 由某条推理规则证成。
如果该推理规则是肯定前件，我们按照 \olref[ded]{thm:deduction-thm} 的证明进行处理。
如果该推理规则是 \QR，我们可知 $!B \ident !C \lif \lforall[x][!D(x)]$，并且形如 $!C \lif !D(a)$ 的公式在此前的推导中出现过，其中 $a$ 不在~$!C$、$!A$ 或 $\Gamma$ 中出现。
于是我们有
\begin{align*}
  \Gamma \cup \{!A\} & \Proves !C \lif !D(a),\\
  \intertext{并且归纳假设适用，即我们有}
    \Gamma & \Proves !A \lif (!C \lif !D(a)).\\
  \intertext{由}
  & \Proves (!A \lif (!C \lif !D(a))) \lif ((!A \land !C) \lif !D(a))\\
  \intertext{以及肯定前件，我们得到}
  \Gamma & \Proves (!A \land !C) \lif !D(a).\\
  \intertext{由于本征变元条件仍然适用，我们可以在该推导中增加一步由 \QR 证成的步骤，从而得到}
    \Gamma & \Proves (!A \land !C) \lif \lforall[x][!D(x)].\\
    \intertext{我们还有}
    & \Proves ((!A \land !C) \lif \lforall[x][!D(x)]) \lif (!A \lif (!C \lif \lforall[x][!D(x)]),\\
    \intertext{因此通过肯定前件，}
    \Gamma & \Proves !A \lif (!C \lif \lforall[x][!D(x)]),
\end{align*}
即 $\Gamma \Proves !A \lif !B$。

我们将 $!B$ 由规则 \QR 证成但形式为 $\lexists[x][!D(x)] \lif !C$ 的情形留作练习。
\end{proof}

\begin{prob}
  完成 \olref[fol][axd][ddq]{thm:deduction-thm-q} 的证明。
\end{prob}

\end{document}